\documentclass[../../../include/open-logic-section]{subfiles}

\begin{document}

\olfileid[hi]{sth}{ordinals}{opps}
\olsection{उत्तरवर्ती और सीमा क्रमसूचक}

अगले कुछ अध्यायों में हम क्रमसूचकों का बहुत उपयोग करेंगे। इसलिए कुछ सरल
धारणाएँ प्रस्तुत करना सहायक होगा।

\begin{defn}
किसी भी क्रमसूचक $\alpha$ का \emph{उत्तरवर्ती}
$\ordsucc{\alpha} =\alpha \cup \{\alpha\}$ है। हम $\alpha$ को
\emph{उत्तरवर्ती} क्रमसूचक कहते हैं यदि किसी क्रमसूचक~$\beta$ के लिए
$\ordsucc{\beta} = \alpha$। हम $\alpha$ को \emph{सीमा} क्रमसूचक तभी और
केवल तभी कहते हैं जब $\alpha$ न रिक्त हो, न उत्तरवर्ती क्रमसूचक।
\end{defn}
\noindent
निम्न परिणाम दिखाता है कि यह \emph{उत्तरवर्ती} की सही धारणा है:

\begin{prop}
किसी भी क्रमसूचक $\alpha$ के लिए:
\begin{enumerate}
	\item $\alpha \in \ordsucc{\alpha}$;
	\item $\ordsucc{\alpha}$ क्रमसूचक है;
	\item ऐसा कोई क्रमसूचक $\beta$ नहीं है कि
	$\alpha \in \beta \in \ordsucc{\alpha}$।
	\end{enumerate}
\end{prop}

\begin{proof}
तुच्छतः $\alpha \in \alpha \cup \{\alpha\} = \ordsucc{\alpha}$। इसी
तरह $\ordsucc{\alpha}$ क्रमसूचकों का संक्रमणशील समुच्चय है और इसलिए
\olref[basic]{corordtransitiveord} से क्रमसूचक है। और
$\alpha \in \beta \in \ordsucc{\alpha}$ असंभव है, क्योंकि तब या तो
$\beta \in \alpha$ अथवा $\beta = \alpha$, जो
\olref[basic]{ordordered} के विरुद्ध है।
\end{proof}
\noindent
यह परिमितेतर आगमन द्वारा प्रमाण के एक रूपांतर की भी अनुमति देता है:

\begin{thm}[सरल परिमितेतर आगमन]\ollabel{simpletransrecursion}
$\phi(x)$ को ऐसा सूत्र मानें कि:
\begin{enumerate}
	\item $\phi(\emptyset)$; और
	\item किसी भी क्रमसूचक $\alpha$ के लिए यदि $\phi(\alpha)$, तो
	$\phi(\ordsucc{\alpha})$; और
	\item यदि $\alpha$ सीमा क्रमसूचक और
	$(\forall \beta \in \alpha)\phi(\beta)$ है, तो $\phi(\alpha)$।
\end{enumerate}
तब $\forall \alpha \phi(\alpha)$।
\end{thm}

\begin{proof}
हम प्रतिधनात्मक सिद्ध करते हैं। मान लें कि कोई क्रमसूचक $\lnot\phi$ है;
$\gamma$ को ऐसा लघुतम क्रमसूचक मानें। तब या तो $\gamma = \emptyset$; अथवा
किसी ऐसे $\alpha$ के लिए $\gamma = \ordsucc{\alpha}$ जिसके लिए
$\phi(\alpha)$; अथवा $\gamma$ सीमा क्रमसूचक है और
$(\forall \beta \in \gamma)\phi(\beta)$।
\end{proof}
\noindent
संकेतन का एक अंतिम अंश आगे सहायक होगा:

\begin{defn}\ollabel{defsupstrict}
यदि $X$ क्रमसूचकों का समुच्चय है, तो
$\supstrict(X) = \bigcup_{\alpha \in X} \ordsucc{\alpha}$।
\end{defn}
\noindent
यहाँ ``lsub'' अंग्रेज़ी ``least strict upper bound'', अर्थात ``लघुतम कठोर
ऊपरी परिबंध'', का संक्षेप है।\footnote{कुछ पुस्तकें इसके लिए
``$\text{sup}(X)$'' का उपयोग करती हैं। किंतु अन्य पुस्तकें
``$\text{sup}(X)$'' का उपयोग लघुतम \emph{अकठोर} ऊपरी परिबंध, अर्थात केवल
$\bigcup X$, के लिए करती हैं। यदि $X$ का महत्तम अवयव $\alpha$ हो, तो ये
धारणाएँ अलग होती हैं: लघुतम \emph{कठोर} ऊपरी परिबंध $\ordsucc{\alpha}$ है,
जबकि लघुतम \emph{अकठोर} ऊपरी परिबंध केवल $\alpha$ है।} निम्न परिणाम इसे
समझाता है:

\begin{prop}
यदि $X$ क्रमसूचकों का समुच्चय है, तो $\supstrict(X)$, $X$ के प्रत्येक
क्रमसूचक से बड़ा लघुतम क्रमसूचक है।
\end{prop}

\begin{proof}
$Y = \Setabs{\ordsucc{\alpha}}{\alpha \in X}$ रखें, जिससे
$\supstrict(X) = \bigcup Y$। चूँकि क्रमसूचक संक्रमणशील हैं और किसी क्रमसूचक
का प्रत्येक सदस्य क्रमसूचक है, इसलिए $\supstrict(X)$ क्रमसूचकों का
संक्रमणशील समुच्चय है और इस प्रकार \olref[basic]{corordtransitiveord} से
क्रमसूचक है।

यदि $\alpha \in X$, तो $\ordsucc{\alpha} \in Y$, अतः
$\ordsucc{\alpha} \subseteq \bigcup Y = \supstrict(X)$ और इसलिए
$\alpha \in \supstrict(X)$। अतः $\supstrict(X)$, $X$ के प्रत्येक
क्रमसूचक से कड़ाई से बड़ा है।

इसके विपरीत यदि $\alpha \in \supstrict(X)$, तो किसी $\beta \in X$ के लिए
$\alpha \in \ordsucc{\beta} \in Y$, जिससे $\alpha \leq \beta \in X$।
अतः $\supstrict(X)$, $X$ का \emph{लघुतम} कठोर ऊपरी परिबंध है।
\end{proof}

\end{document}
